%! AP Statistics — Unit 1, Lesson 1.2 — Homework
\documentclass[10pt]{article}
\usepackage{apstats-article}
\usepackage{apstats-boxes}

%=============================================================================
\begin{document}

\pageheader{Unit 1, Lesson 1.2}{Homework: The Language of Variation --- Variables}
\namedateperiod

\vspace{0.1in}

\begin{tcolorbox}[colback=sky, colframe=navy, boxrule=0.8pt, arc=2mm,
    left=3mm, right=3mm, top=1.5mm, bottom=1.5mm]
\small For each problem, classify the variable(s) and justify your reasoning.
A correct type with no justification earns partial credit only ---
AP free-response always requires explanation.
\end{tcolorbox}

\vspace{0.14in}

% ════════════════════════════════════════════════════════════════════════════
% PART A: NBA DATASET
% ════════════════════════════════════════════════════════════════════════════
\begin{blurbbox}[Part A \quad NBA Player Dataset]{navylight}

\small The table below shows data for five NBA players from a recent season.

\vspace{8pt}
\renewcommand{\arraystretch}{1.5}
\begin{tabularx}{\linewidth}{@{} >{\bfseries}p{0.16\linewidth}
                                    C{0.13\linewidth} C{0.13\linewidth}
                                    C{0.16\linewidth} C{0.13\linewidth}
                                    C{0.14\linewidth} @{}}
\rowcolor{sky}
\textbf{Name} & \textbf{Team} & \textbf{Pts/game} &
\textbf{Position} & \textbf{Games} & \textbf{Jersey \#} \\
\midrule
Player A & Lakers  & 27.4 & Forward & 68 & 23 \\
Player B & Celtics & 19.8 & Guard   & 72 & 11 \\
Player C & Warriors& 31.2 & Guard   & 65 & 30 \\
Player D & Heat    & 14.5 & Center  & 79 & 3  \\
Player E & Bucks   & 23.7 & Forward & 71 & 34 \\
\end{tabularx}

\vspace{12pt}

Classify each of the six variables (Name, Team, Pts/game, Position, Games played,
Jersey \#). For each, write the type \textbf{and} a one-phrase justification.

\vspace{8pt}
\renewcommand{\arraystretch}{2.2}
\begin{tabularx}{\linewidth}{@{} >{\bfseries}p{0.2\linewidth}
                                    C{0.14\linewidth} X @{}}
\rowcolor{sky}
\textbf{Variable} & \textbf{Type} & \textbf{Justification} \\
\midrule
Name        & \blank{2cm} & \blank{8.5cm} \\
Team        & \blank{2cm} & \blank{8.5cm} \\
Pts/game    & \blank{2cm} & \blank{8.5cm} \\
Position    & \blank{2cm} & \blank{8.5cm} \\
Games played& \blank{2cm} & \blank{8.5cm} \\
Jersey \#   & \blank{2cm} & \blank{8.5cm} \\
\end{tabularx}

\vspace{6pt}
\textbf{Which variable is most likely to cause confusion? Explain why.}\\[4pt]
\writeline\\[1pt]\writeline

\end{blurbbox}

\vspace{0.14in}

% ════════════════════════════════════════════════════════════════════════════
% PART B: AP-STYLE MULTIPLE CHOICE
% ════════════════════════════════════════════════════════════════════════════
\begin{blurbbox}[Part B \quad AP-Style Multiple Choice]{greenacc}
\small Circle the letter of the best answer. Write one sentence explaining your choice.

\vspace{10pt}
\textbf{Question 1.}\enspace
A school survey records each student's grade level (9, 10, 11, or 12).
Which of the following is correct?

\begin{enumerate}[label=\textbf{(\Alph*)}, leftmargin=2em, itemsep=4pt]
  \item Grade level is quantitative continuous because it is recorded as a number.
  \item Grade level is quantitative discrete because grades are countable whole numbers.
  \item Grade level is categorical because the numbers label groups, not measured amounts.
  \item Grade level cannot be classified without more information.
\end{enumerate}
\textcolor{slate}{\scriptsize Explanation:}\;\writeline

\vspace{12pt}
\textbf{Question 2.}\enspace
A meteorologist records the \textit{type of precipitation} (rain, snow, sleet, none) and
the \textit{total precipitation amount} (in inches) for each day in January.
Which statement correctly classifies both variables?

\begin{enumerate}[label=\textbf{(\Alph*)}, leftmargin=2em, itemsep=4pt]
  \item Both variables are categorical.
  \item Type of precipitation is categorical; total amount is quantitative continuous.
  \item Type of precipitation is quantitative discrete; total amount is quantitative continuous.
  \item Both variables are quantitative.
\end{enumerate}
\textcolor{slate}{\scriptsize Explanation:}\;\writeline

\vspace{12pt}
\textbf{Question 3.}\enspace
A doctor records a patient's \textit{pain level} on a scale of 0--10 and the
\textit{number of days since surgery}. A classmate says both are quantitative discrete.
Is the classmate correct?

\begin{enumerate}[label=\textbf{(\Alph*)}, leftmargin=2em, itemsep=4pt]
  \item Yes --- both are whole numbers, so both are quantitative discrete.
  \item No --- pain level is categorical (ordinal); days since surgery is quantitative discrete.
  \item No --- pain level is debatable (categorical or QD); days since surgery is quantitative discrete.
  \item No --- both are quantitative continuous because they are measured.
\end{enumerate}
\textcolor{slate}{\scriptsize Explanation:}\;\writeline

\end{blurbbox}

\vspace{0.14in}

% ════════════════════════════════════════════════════════════════════════════
% EXTENSION
% ════════════════════════════════════════════════════════════════════════════
\begin{extensionbox}
\small
\textbf{Optional Extension --- Real Data from the CDC}\\[4pt]
Go to \texttt{wonder.cdc.gov} or \texttt{data.census.gov} and find any publicly available dataset.
Identify \textbf{three categorical} and \textbf{three quantitative} variables.
For the quantitative variables, note whether each is discrete or continuous.

\vspace{8pt}
\renewcommand{\arraystretch}{2.0}
\begin{tabularx}{\linewidth}{@{} >{\bfseries}p{0.14\linewidth}
                                    >{\raggedright\arraybackslash}p{0.36\linewidth}
                                    X @{}}
\rowcolor{sky}
& \textbf{Variable name} & \textbf{Notes / justification} \\
\midrule
Cat. 1 & \blank{7.5cm} & \blank{3.5cm} \\
Cat. 2 & \blank{7.5cm} & \blank{3.5cm} \\
Cat. 3 & \blank{7.5cm} & \blank{3.5cm} \\
Quant. 1 & \blank{7.5cm} & \blank{3.5cm} \\
Quant. 2 & \blank{7.5cm} & \blank{3.5cm} \\
Quant. 3 & \blank{7.5cm} & \blank{3.5cm} \\
\end{tabularx}

\vspace{6pt}
\textcolor{slate}{\scriptsize Source / URL:\;\blank{11cm}}

\vspace{4pt}
\textcolor{slate}{\scriptsize We will return to categorical variables in Lesson 1.3 and quantitative variables in Lesson 1.5.}
\end{extensionbox}

\end{document}
